\documentclass{article}
\usepackage[margin=1in]{geometry}
\usepackage{amsmath,amssymb}
\usepackage[most]{tcolorbox}
\usepackage[hidelinks]{hyperref}

\newcommand{\BookName}{Pattern Recognition and Machine Learning}
\newcommand{\BookAuthor}{Christopher M. Bishop}
\newcommand{\ChapterName}{Introduction}
\newcommand{\ExerciseID}{1.9}

\title{Chapter: \ChapterName}
\author{Ahonakpon Guy Gbaguidi}
\date{}

\begin{document}

\maketitle

\begin{center}
  \large\textbf{Exercise \ExerciseID}

  \vspace{0.5em}

  \normalsize\BookName\ - \BookAuthor
\end{center}

\begin{tcolorbox}[breakable,colback=gray!5,colframe=gray!60,title=Solution]

\subsubsection*{Part (a)}
Let's show that the mode (i.e. the maximum) of the Gaussian distribution given by (1.46) is \( \mu \).

Let's restate the equation (1.46) for clarity:

\begin{equation}
  \mathcal{N}(x|\mu,\sigma^2) = \frac{1}{\left(2\pi\sigma^2\right)^{1/2}} \exp\left\{-\frac{1}{2\sigma^2}(x-\mu)^2\right\} \tag{1.46}
\end{equation}

To find the mode of the Gaussian distribution, we need to find the value of \( x \) that maximizes 
\( \mathcal{N}(x|\mu,\sigma^2) \). Since the exponential function is monotonically increasing, we 
can focus on maximizing the exponent:
\begin{equation}
  -\frac{1}{2\sigma^2}(x-\mu)^2
\end{equation}

To maximize this expression, we need to minimize \( (x-\mu)^2 \) because \(-\frac{1}{2\sigma^2} < 0\). 
The minimum value of \( (x-\mu)^2 \) is 0 because \((x-\mu)^2 \geq 0\) for all \( x \), 
which occurs when \( x = \mu \). Therefore, the mode of the Gaussian distribution is \( \mu \) 
and it is unique because \(\sigma^2 > 0\).

\vspace{1em}

\subsubsection*{Part (b)}
Similarly, we can show that the mode of the multivariate Gaussian distribution given by (1.52) is \( \boldsymbol{\mu} \).
Let's restate the equation (1.52) for clarity:

\begin{equation}
  \mathcal{N}(\mathbf{x}|\boldsymbol{\mu},\Sigma) = \frac{1}{(2\pi)^{D/2} |\Sigma|^{1/2}} \exp\left\{-\frac{1}{2}(\mathbf{x}-\boldsymbol{\mu})^T \Sigma^{-1} (\mathbf{x}-\boldsymbol{\mu})\right\} \tag{1.52}
\end{equation}

where \( \mathbf{x} \) is a \( D \)-dimensional vector, \( \boldsymbol{\mu} \) is the mean vector, and \( \Sigma \) is the covariance \( D \times D \) matrix and \( |\Sigma| \) is the determinant of \( \Sigma \).

\vspace{0.5em}
To find the mode of the multivariate Gaussian distribution, we need to find the value of \( \mathbf{x} \) 
that maximizes \( \mathcal{N}(\mathbf{x}|\boldsymbol{\mu},\Sigma) \). 
Similar to the univariate case, we can focus on maximizing the exponent:
\begin{equation}
  -\frac{1}{2}(\mathbf{x}-\boldsymbol{\mu})^T \Sigma^{-1} (\mathbf{x}-\boldsymbol{\mu})
\end{equation}

To maximize this expression, we need to minimize 
\( (\mathbf{x}-\boldsymbol{\mu})^T \Sigma^{-1} (\mathbf{x}-\boldsymbol{\mu}) \).
Since \( \Sigma \) is a positive definite matrix, \( \Sigma^{-1} \) is also positive definite, which means 
that \( (\mathbf{x}-\boldsymbol{\mu})^T \Sigma^{-1} (\mathbf{x}-\boldsymbol{\mu}) \geq 0 \) for all 
\( \mathbf{x} \). The minimum value of this expression is 0, which occurs when \( \mathbf{x} = \boldsymbol{\mu} \). 
Therefore, the mode of the multivariate Gaussian distribution is \( \boldsymbol{\mu} \) and it is unique 
because \( \Sigma \) is positive definite.
\vspace{1em}
\paragraph{Justification.}
Since \(\Sigma\) is symmetric positive definite there exists an invertible symmetric square root \(\Sigma^{1/2}\) with
\(\Sigma=\Sigma^{1/2}\Sigma^{1/2}\) and \(\Sigma^{-1}=\Sigma^{-1/2}\Sigma^{-1/2}\). Hence
\[
(\mathbf{x}-\boldsymbol{\mu})^T\Sigma^{-1}(\mathbf{x}-\boldsymbol{\mu})
= \|\Sigma^{-1/2}(\mathbf{x}-\boldsymbol{\mu})\|^2 \ge 0.
\]
Equality holds iff \(\Sigma^{-1/2}(\mathbf{x}-\boldsymbol{\mu})=\mathbf{0}\), and since \(\Sigma^{-1/2}\) is invertible this is equivalent to
\(\mathbf{x}=\boldsymbol{\mu}\).

\paragraph{Resources to learn more}
\begin{itemize}
  \item \href{https://en.wikipedia.org/wiki/Positive-definite_matrix#Square_root_of_a_positive-definite_matrix}{Square root of a positive definite matrix}
  \item \href{https://en.wikipedia.org/wiki/Positive-definite_matrix#Cholesky_decomposition}{Cholesky decomposition}
\end{itemize}

\end{tcolorbox}

\end{document}
